\documentclass[../../../../main.tex]{subfiles}
\begin{document}

\begin{flushright}
\footnotesize\textit{【自動文字起こし・要確認】}
\end{flushright}

[解] $k \ge 0$である.
(1) $F(x) = f(x) - x$ とおくと, $F'(x) = f'(x) - 1 < 0 \ (\because (\text{ハ}))$ から, $F(x)$は
単調減少. これと $F(a) = f(a) - a \ge a - a = 0, \ F(b) = f(b) - b \le 0$
から, $F(x) = 0$をみたす$x$が $[a,b]$の中に唯1つ存在し, これが
$C_0$に他ならない. (終)

(2) 平均値の定理から $d_n \neq C_0$のとき,
$$ f(d_n) - f(C_0) = f'(p) (d_n - C_0) $$
をみたす $f'(p)$がある. (ハ)から $|f'(p)| \le k$ だから, 両辺
絶対値をとってセイリして,
$$ |d_{n+1} - C_0| \le k |d_n - C_0| \quad (\because f(d_n) = d_{n+1}, f(C_0) = C_0) $$
くり返し用いて
$$ |d_n - C_0| \le k^n |d_0 - C_0| $$
又, $d_n = C_0$の時, 明らかに成立.
以上から示された. (終)

\end{document}
